\chapter{系统需求分析与总体设计}

本章在前期技术研究的基础上，对基于鸿蒙系统与MindSpore的智能图像分类应用进行需求分析，明确系统的功能需求和非功能需求，并在此基础上进行系统总体架构设计、技术路线设计和模块划分，为后续的详细设计与实现奠定基础。

\section{系统需求分析}

\subsection{功能需求分析}

根据智能图像分类应用的实际使用场景和用户需求，本系统需具备以下功能：

（1）图像选择功能。用户应能够从设备相册中选取待分类的图片，也可以通过设备摄像头实时拍摄图片进行分类。系统需支持常见图片格式（JPEG、PNG等）的导入，并提供图片预览功能。

（2）图像预处理功能。系统应能对用户选取的原始图片进行预处理操作，包括图像缩放、裁剪、归一化等，使输入图像满足MindSpore Lite模型的推理要求。

（3）图像分类功能。系统核心功能，应能调用MindSpore Lite推理引擎对预处理后的图像进行前向推理，输出图像所属类别的概率分布和预测标签。

（4）结果展示功能。系统应能以直观的方式展示分类结果，包括预测类别、置信度（Top-1和Top-5），支持以列表或图表形式呈现各类别的概率分布。

（5）历史记录功能。系统应能自动保存每次分类的记录，包括原始图片、分类结果、置信度和分类时间等信息，支持历史记录的浏览、搜索和删除操作。

\subsection{非功能需求分析}

除上述功能需求外，系统还应满足以下非功能需求：

（1）实时性要求。图像分类推理响应时间应小于100ms（不含图像选择时间），确保用户获得流畅的交互体验。

（2）准确性要求。在CIFAR-10测试集上，分类准确率应达到90\%以上，满足实际应用的基本精度需求。

（3）资源占用要求。应用运行时的内存占用应控制在200MB以内，避免对设备其他功能造成影响。

（4）离线可用性要求。应用应支持完全离线运行，在无网络连接环境下仍能正常完成图像分类功能，所有模型推理均在端侧完成。

（5）兼容性要求。应用应兼容鸿蒙系统4.0及以上版本，支持不同分辨率的鸿蒙设备。

（6）稳定性要求。应用应能稳定运行，连续执行100次分类任务不出现崩溃或内存泄漏。

\section{系统总体架构设计}

基于上述需求分析，本系统采用三层架构进行设计，分别为表示层、业务逻辑层和数据层。系统总体架构如图\ref{fig:architecture}所示。

\begin{figure}[htbp]
	\centering
	\begin{tikzpicture}[
		box/.style={draw, rounded corners, minimum width=3.2cm, minimum height=1cm, align=center, font=\small},
		layerbox/.style={draw, rounded corners, minimum width=13cm, minimum height=2.2cm, align=center},
		arr/.style={->, thick, >=stealth}
	]

	\def\layerw{13}
	\def\boxw{2.8}

	\node[layerbox, fill=blue!10] (pres) at (0, 4.5) {};
	\node[font=\bfseries] at (-4.8, 5.2) {表示层};
	\node[box, fill=white] (ui1) at (-3.5, 4.5) {图像选择\\界面};
	\node[box, fill=white] (ui2) at (0, 4.5) {分类结果\\展示界面};
	\node[box, fill=white] (ui3) at (3.5, 4.5) {历史记录\\管理界面};

	\node[layerbox, fill=green!10] (biz) at (0, 1.5) {};
	\node[font=\bfseries] at (-4.8, 2.2) {业务逻辑层};
	\node[box, fill=white] (b1) at (-3.5, 1.5) {图像预处理\\模块};
	\node[box, fill=white] (b2) at (0, 1.5) {MindSpore Lite\\推理引擎};
	\node[box, fill=white] (b3) at (3.5, 1.5) {NAPI\\桥接层};

	\node[layerbox, fill=orange!10] (data) at (0, -1.5) {};
	\node[font=\bfseries] at (-4.8, -0.8) {数据层};
	\node[box, fill=white] (d1) at (-3.5, -1.5) {模型文件\\存储};
	\node[box, fill=white] (d2) at (0, -1.5) {历史记录\\数据库};
	\node[box, fill=white] (d3) at (3.5, -1.5) {配置文件\\管理};

	\draw[arr] (ui1) -- (b1);
	\draw[arr] (ui2) -- (b2);
	\draw[arr] (ui3) -- (b2);
	\draw[arr] (b1) -- (b2);
	\draw[arr] (b2) -- (b3);
	\draw[arr] (d1) -- (b3);
	\draw[arr] (d2) -- (ui3);
	\draw[arr] (d3) -- (b3);

	\end{tikzpicture}
	\caption{系统总体架构图}
	\label{fig:architecture}
\end{figure}

\subsection{表示层}

表示层基于鸿蒙ArkUI声明式开发框架构建，负责用户界面的展示和交互逻辑的处理。表示层包含以下三个核心界面组件：

（1）图像选择界面。提供相册选图和相机拍摄两种图像获取方式，支持图片格式过滤和预览功能，采用Grid布局展示相册缩略图。

（2）分类结果展示界面。以卡片式布局展示分类结果，包括预测类别标签、Top-1/Top-5置信度数值和概率分布条形图，提供结果分享和保存操作入口。

（3）历史记录管理界面。采用List列表组件展示历史分类记录，支持按时间排序、关键词搜索和批量删除操作。

\subsection{业务逻辑层}

业务逻辑层是系统的核心层，负责图像处理和模型推理两大核心业务。业务逻辑层包含以下三个关键模块：

（1）图像预处理模块。对用户选取的原始图像进行尺寸缩放（统一至224$\times$224像素）、中心裁剪、颜色通道转换和BGR-RGB转换，以及像素值归一化处理，将处理后的数据转换为模型所需的Tensor格式。

（2）MindSpore Lite推理引擎。加载端侧优化后的ResNet-50量化模型，执行前向推理计算，输出10个类别的概率分布向量。推理引擎通过鸿蒙NAPI接口与表示层进行数据交互。

（3）NAPI桥接层。作为ArkTS与C++之间的通信桥梁，负责将ArkTS侧的图像数据和处理请求传递给C++侧的推理引擎，同时将推理结果返回给ArkTS侧。NAPI桥接层实现了跨语言调用的封装。

\subsection{数据层}

数据层负责系统数据的持久化存储和管理，包含以下三个存储组件：

（1）模型文件存储。采用鸿蒙应用的rawfile目录存储MindSpore Lite格式的模型文件（.ms格式），应用启动时从rawfile目录将模型文件复制到应用沙箱路径下供推理引擎加载。

（2）历史记录数据库。基于鸿蒙关系型数据库（RelationalStore）实现历史分类记录的增删改查操作，存储字段包括图片路径、预测类别、置信度、分类时间等。

（3）配置文件管理。采用JSON格式存储应用配置信息，包括模型版本号、分类类别标签映射表、预处理参数（均值、标准差）等。

\section{技术路线设计}

本系统的技术路线主要围绕模型训练、模型优化、端侧部署和应用开发四个环节展开，各环节所采用的关键技术如下：

（1）模型训练——MindSpore框架。采用MindSpore深度学习框架搭建ResNet-50卷积神经网络，在CIFAR-10数据集上进行监督学习训练。MindSpore提供了自动微分、自动并行等特性，支持灵活的训练策略配置，能够高效完成模型训练过程。训练阶段采用交叉熵损失函数和Adam优化器，学习率采用余弦退火策略进行动态调整。

（2）模型推理——MindSpore Lite。将训练完成的模型通过MindSpore Lite转换工具转换为.ms格式文件，并采用训练后量化（Post-Training Quantization）策略将模型权重从FP32量化为INT8，在保持模型精度下降不超过1\%的前提下，将模型文件大小压缩至原始大小的1/4，推理速度提升约2倍。

（3）应用开发——ArkTS语言与ArkUI框架。采用鸿蒙官方推荐的ArkTS语言进行应用开发，使用ArkUI声明式开发范式构建用户界面。ArkTS是在TypeScript基础上扩展的方言，具有强类型检查和良好的开发体验。ArkUI提供丰富的组件库，支持Flex、Grid、List等多种布局方式，适合构建响应式界面。

（4）跨语言调用——NAPI接口。MindSpore Lite推理引擎以C++动态链接库形式提供，需要通过鸿蒙NAPI（Native API）机制实现ArkTS与C++之间的互操作。NAPI提供了类型转换、异常处理和异步回调等能力，是鸿蒙系统中连接ArkTS层与Native层的关键桥梁。本系统通过NAPI封装了模型初始化、推理执行和模型释放三个核心接口。

\section{系统模块划分}

根据系统的功能需求和架构设计，将系统划分为以下五个功能模块。各模块的功能描述如表\ref{tab:modules}所示。

\begin{table}[htbp]
	\centering
	\caption{系统模块划分及功能描述}
	\label{tab:modules}
	\begin{tabular}{p{3cm}p{9cm}}
		\toprule
		模块名称 & 功能描述 \\
		\midrule
		图像获取模块 & 负责图像的获取功能，包括从设备相册选择图片和通过摄像头实时拍摄图片两种方式。调用鸿蒙PhotoViewPicker组件实现相册选图，调用Camera组件实现拍照功能，并将获取的图像统一转换为PixelMap格式供后续处理。 \\
		\addlinespace
		图像预处理模块 & 负责对原始图像进行预处理，包括图像缩放（224$\times$224）、中心裁剪、颜色通道转换（BGR转RGB）和像素值归一化（均值为[0.485, 0.456, 0.406]，标准差为[0.229, 0.224, 0.225]）。将预处理后的数据组装为Float32类型的Tensor输入。 \\
		\addlinespace
		模型推理模块 & 负责模型加载和推理执行。模型推理模块通过NAPI接口调用MindSpore Lite C++ API，完成模型上下文创建、会话配置、Tensor绑定的初始化流程，执行前向推理计算，并将输出Tensor解码为类别概率分布。 \\
		\addlinespace
		结果展示模块 & 负责分类结果的可视化展示，包括预测类别标签、Top-1和Top-5置信度数值、各类别概率分布条形图。支持结果分享和一键保存至历史记录功能。 \\
		\addlinespace
		历史管理模块 & 负责分类历史记录的持久化存储和管理。基于鸿蒙RelationalStore关系型数据库实现历史记录的增删改查操作，存储字段包括图片缩略图路径、预测类别、置信度、分类时间戳。支持按时间倒序排列和关键词搜索。 \\
		\bottomrule
	\end{tabular}
\end{table}

各模块之间的交互关系如下：用户通过图像获取模块选取或拍摄图片后，图像数据传入图像预处理模块进行标准化处理；预处理后的Tensor数据由模型推理模块执行推理计算，输出分类概率分布；结果展示模块将推理结果以可视化方式呈现给用户，同时将结果数据传递给历史管理模块进行持久化存储。整个数据流转过程遵循单向数据流原则，保证了系统的可维护性和可扩展性。